\documentclass[tikz,border=5pt]{standalone}
\usepackage{ctex}
\usepackage{amsmath}
\usetikzlibrary{arrows.meta,calc,positioning}

\begin{document}
\begin{tikzpicture}[
    >=Stealth,
    line join=round,
    line cap=round,
    font=\small
]

% =========================
% 参数
% =========================
\def\L{5.2}      % 船长
\def\W{2.0}      % 船宽
\def\ax{2.0}     % 艇体坐标轴长度
\def\thrX{-1.6}  % 推进器x位置
\def\thrY{0.75}  % 推进器y偏置
\def\arr{1.8}    % 推力箭头长度

% =========================
% 船体轮廓（俯视）
% =========================
\draw[thick]
    (-0.45*\L,-0.5*\W)
    -- (0.22*\L,-0.5*\W)
    -- (0.48*\L,-0.28*\W)
    -- (0.62*\L,0)
    -- (0.48*\L,0.28*\W)
    -- (0.22*\L,0.5*\W)
    -- (-0.45*\L,0.5*\W)
    -- cycle;

% 中轴虚线
\draw[dashed] (-0.42*\L,0) -- (0.56*\L,0);

% =========================
% 质心 G
% =========================
\coordinate (G) at (0,0);
\fill (G) circle (1.2pt);
\node[above right=1pt] at (G) {$G$};

% =========================
% 艇体坐标系
% =========================
\draw[->, thick] (G) -- (\ax,0) node[above] {$x_b$};
\draw[->, thick] (G) -- (0,-\ax) node[right] {$y_b$};

% =========================
% 左右推进器
% =========================
\coordinate (TLp) at (\thrX,\thrY);
\coordinate (TRp) at (\thrX,-\thrY);

% 推进器本体
\draw[thick, fill=white] ($(TLp)+(-0.18,0.12)$) rectangle ($(TLp)+(0.18,-0.12)$);
\draw[thick, fill=white] ($(TRp)+(-0.18,0.12)$) rectangle ($(TRp)+(0.18,-0.12)$);

\node[left=2pt] at (TLp) {L};
\node[left=2pt] at (TRp) {R};

% 推力箭头
\draw[->, thick] (TLp) -- ++(\arr,0) node[above] {$T_L$};
\draw[->, thick] (TRp) -- ++(\arr,0) node[below] {$T_R$};

% =========================
% 推进器到中轴线的距离 b
% =========================
\draw[dashed] (TLp) -- (\thrX,0);
\draw[dashed] (TRp) -- (\thrX,0);

\draw[<->] (\thrX+0.45,0) -- (\thrX+0.45,\thrY)
    node[midway,right] {$b$};
\draw[<->] (\thrX+0.45,0) -- (\thrX+0.45,-\thrY)
    node[midway,right] {$b$};

% =========================
% 合力 X
% =========================
\draw[->, very thick] (0.95,1.45) -- ++(1.6,0) node[above] {$X$};

% =========================
% 偏航力矩 N（示意）
% =========================
\draw[->, thick] (1.4,-0.2) arc[start angle=-20,end angle=250,radius=0.7];
\node[right] at (1.9,0.65) {$N$};

% =========================
% 说明性辅助标注（可选）
% =========================
\node[below] at (0,-1.45) {双推进器差速输入示意};

\end{tikzpicture}
\end{document}